\subsection{Perkalian Matriks}

Operasi matriks penting berikutnya yang akan kita pelajari adalah perkalian
matriks. Operasi
perkalian matriks merupakan salah satu operasi matriks yang paling penting dan
bermanfaat. Di sepanjang bagian ini, kita juga akan
menunjukkan hubungan perkalian matriks dengan sistem
persamaan linear.

Pertama, kita memberikan definisi formal vektor baris dan vektor kolom. 

\begin{definition}{Vektor Baris dan Vektor Kolom}\label{def:rowandcolumnvectors}
Matriks berukuran $n\times 1$ atau $1\times n$ disebut \textbf{vektor}. Jika $\x$ adalah matriks semacam itu, kita menulis $x_{i}$ untuk
menyatakan entri $\x$ pada baris $i$ dari matriks kolom, atau kolom $i$ dari matriks baris.
\index{vektor}
\index{vektor!kolom} 

Matriks $n\times 1$
\begin{equation*}
\x=\leftB
\begin{array}{c}
x_{1} \\
\vdots \\
x_{n}
\end{array}
\rightB
\end{equation*}
disebut
\index{vektor!vektor baris} \textbf{vektor kolom.} Matriks $1\times n$
\begin{equation*}
\x = \leftB
\begin{array}{ccc}
x_{1} & \cdots & x_{n}
\end{array}
\rightB
\end{equation*}
disebut \textbf{vektor baris}.
\end{definition}

Di seluruh teks ini, kita dapat menggunakan istilah \textbf{vektor} saja untuk menyebut vektor kolom maupun vektor baris.
Jika demikian, konteks akan memperjelas jenis vektor yang dimaksud.

Dalam bab ini, kita kembali menggunakan konsep kombinasi linear
vektor seperti dalam Definisi \ref{def:linearcombination}. Dalam konteks ini, kombinasi linear adalah suatu jumlah
yang terdiri atas vektor-vektor yang dikalikan dengan skalar. Sebagai contoh,
\begin{equation*}
\leftB
\begin{array}{r}
50 \\
122
\end{array}
\rightB
=
7\leftB
\begin{array}{r}
1 \\
4
\end{array}
\rightB +8\leftB
\begin{array}{r}
2 \\
5
\end{array}
\rightB +9\leftB
\begin{array}{r}
3 \\
6
\end{array}
\rightB
\end{equation*}
adalah kombinasi linear dari tiga vektor. 

Ternyata, kita dapat menyatakan sistem persamaan linear apa pun sebagai kombinasi linear vektor. Bahkan,
vektor-vektor yang akan kita gunakan hanyalah kolom-kolom matriks diperbesar yang bersesuaian! 

\begin{definition}{Bentuk Vektor Sistem Persamaan Linear}\label{def:vectorform}
Misalkan kita memiliki sistem persamaan yang diberikan oleh
\begin{equation*}
\begin{array}{c}
a_{11}x_{1}+\cdots +a_{1n}x_{n}=b_{1} \\
\vdots \\
a_{m1}x_{1}+\cdots +a_{mn}x_{n}=b_{m}
\end{array}
\end{equation*}
Kita dapat menyatakan sistem ini dalam \textbf{bentuk vektor} \index{sistem persamaan!bentuk vektor} sebagai berikut:
\begin{equation*}
x_1
\leftB
\begin{array}{c}
a_{11}\\
a_{21}\\
\vdots \\
a_{m1}
\end{array}
\rightB
+
x_2
\leftB
\begin{array}{c}
a_{12}\\
a_{22}\\
\vdots \\
a_{m2}
\end{array}
\rightB
+
\cdots
+
x_n
\leftB
\begin{array}{c}
a_{1n}\\
a_{2n}\\
\vdots \\
a_{mn}
\end{array}
\rightB
=
\leftB
\begin{array}{c}
b_1\\
b_2\\
\vdots \\
b_m
\end{array}
\rightB
\end{equation*}
\end{definition}

Perhatikan bahwa setiap vektor yang digunakan di sini merupakan satu kolom dari matriks diperbesar yang bersesuaian. Ada satu vektor untuk setiap variabel dalam sistem,
beserta vektor konstanta. 

Bentuk penting pertama dari perkalian matriks adalah mengalikan matriks dengan vektor.
Perhatikan hasil kali
\begin{equation*}
\leftB
\begin{array}{rrr}
1 & 2 & 3 \\
4 & 5 & 6
\end{array}
\rightB \leftB
\begin{array}{r}
7 \\
8 \\
9
\end{array}
\rightB
\end{equation*}
Kita akan segera melihat bahwa hasil ini sama dengan
\begin{equation*}
7\leftB
\begin{array}{c}
1 \\
4
\end{array}
\rightB +8\leftB
\begin{array}{c}
2 \\
5
\end{array}
\rightB +9\leftB
\begin{array}{c}
3 \\
6
\end{array}
\rightB =\leftB
\begin{array}{c}
50 \\
122
\end{array}
\rightB
\end{equation*}

Secara umum,
\begin{eqnarray*}
\leftB
\begin{array}{ccc}
a_{11} & a_{12} & a_{13} \\
a_{21} & a_{22} & a_{23}
\end{array}
\rightB \leftB
\begin{array}{c}
x_{1} \\
x_{2} \\
x_{3}
\end{array}
\rightB &=&\allowbreak x_{1}\leftB
\begin{array}{c}
a_{11} \\
a_{21}
\end{array}
\rightB +x_{2}\leftB
\begin{array}{c}
a_{12} \\
a_{22}
\end{array}
\rightB +x_{3}\leftB
\begin{array}{c}
a_{13} \\
a_{23}
\end{array}
\rightB \\
&=&\leftB
\begin{array}{c}
a_{11}x_{1}+a_{12}x_{2}+a_{13}x_{3} \\
a_{21}x_{1}+a_{22}x_{2}+a_{23}x_{3}
\end{array}
\rightB 
\end{eqnarray*}
Jadi, Anda mengambil $x_{1}$ kali kolom pertama, menambahkannya dengan $x_{2}$ kali
kolom kedua, lalu akhirnya menambahkan $x_{3}$ kali kolom ketiga. Jumlah di atas adalah kombinasi linear dari kolom-kolom matriks tersebut.
Ketika Anda mengalikan sebuah matriks di sebelah kiri dengan sebuah vektor di sebelah kanan,
bilangan-bilangan pembentuk vektor itu merupakan skalar-skalar yang digunakan dalam
kombinasi linear kolom seperti diilustrasikan di atas.

Berikut adalah definisi formal cara mengalikan matriks $m\times
n $ dengan vektor kolom $ n\times 1 $.

\begin{definition}{Perkalian Matriks dengan Vektor}\label{def:multiplicationvectormatrix}
Misalkan $A=\leftB a_{ij} \rightB$ adalah matriks $m\times n$ dan $\x$
adalah matriks $n\times 1$ yang diberikan oleh
\begin{equation*}
A=\leftB A_{1} \cdots A_{n}\rightB, \quad \x = \leftB
\begin{array}{r}
x_{1} \\
\vdots \\
x_{n}
\end{array}
\rightB 
\end{equation*}

Maka hasil kali $A\x$ adalah vektor kolom $m\times 1$
\index{perkalian matriks!vektor} yang sama dengan kombinasi linear
kolom-kolom $A$ berikut:
\begin{equation*}
x_{1}A_{1}+x_{2}A_{2}+\cdots +x_{n}A_{n} = 
\sum_{j=1}^{n}x_{j}A_{j}  
\end{equation*}
\end{definition}

Jika kita menuliskan kolom-kolom $A$ berdasarkan entri-entrinya, kolom-kolom tersebut berbentuk
\begin{equation*}
A_{j}  =
\leftB
\begin{array}{c}
a_{1j} \\
a_{2j} \\
\vdots \\
a_{mj}
\end{array}
\rightB 
\end{equation*}
Maka kita dapat menuliskan hasil kali $A\x$ sebagai
\begin{equation*}
A\x = 
x_{1}\leftB
\begin{array}{c}
a_{11} \\
a_{21} \\
\vdots \\
a_{m1}
\end{array}
\rightB + x_{2}\leftB
\begin{array}{c}
a_{12} \\
a_{22} \\
\vdots \\
a_{m2}
\end{array}
\rightB +\cdots + x_{n}\leftB
\begin{array}{c}
a_{1n} \\
a_{2n} \\
\vdots \\
a_{mn}
\end{array}
\rightB 
\end{equation*}

Perhatikan bahwa perkalian matriks $m \times n$ dengan vektor $n \times 1$ menghasilkan vektor $m \times 1$.

Berikut sebuah contoh.

\begin{example}{Perkalian Matriks dengan Vektor}\label{exa:vectormultbymatrix}
Hitung hasil kali $A\x$ untuk
\begin{equation*}
A = \leftB
\begin{array}{rrrr}
1 & 2 & 1 & 3 \\
0 & 2 & 1 & -2 \\
2 & 1 & 4 & 1
\end{array}
\rightB,\,  \x =  \leftB
\begin{array}{r}
1 \\
2 \\
0 \\
1
\end{array}
\rightB 
\end{equation*}
\end{example}

\begin{solution} Kita akan menggunakan Definisi \ref{def:multiplicationvectormatrix} untuk menghitung hasil kali tersebut.
Oleh karena itu, kita menghitung hasil kali $A\x$ sebagai berikut. 
\begin{eqnarray*}
& 1\leftB
\begin{array}{r}
1 \\
0 \\
2
\end{array}
\rightB + 2\leftB
\begin{array}{r}
2 \\
2 \\
1
\end{array}
\rightB + 0\leftB
\begin{array}{r}
1 \\
1 \\
4
\end{array}
\rightB + 1 \leftB
\begin{array}{r}
3 \\
-2\\
1
\end{array} 
\rightB \\ 
&=
\leftB
\begin{array}{r}
1 \\
0 \\
2
\end{array}
\rightB + \leftB
\begin{array}{r}
4 \\
4 \\
2
\end{array}
\rightB + \leftB
\begin{array}{r}
0 \\
0 \\
0
\end{array}
\rightB +  \leftB
\begin{array}{r}
3 \\
-2\\
1
\end{array} 
\rightB \\
&=
\leftB
\begin{array}{r}
8 \\
2 \\
5
\end{array}
\rightB
\end{eqnarray*}
\end{solution}

Dengan menggunakan operasi di atas, kita juga dapat menuliskan sistem persamaan linear dalam \textbf{bentuk matriks}. Dalam bentuk ini,
kita menyatakan sistem sebagai matriks yang dikalikan dengan vektor. Perhatikan definisi berikut.

\begin{definition}{Bentuk Matriks Sistem Persamaan Linear}\label{def:matrixform}
Misalkan kita memiliki sistem persamaan yang diberikan oleh
\begin{equation*}
\begin{array}{c}
a_{11}x_{1}+\cdots +a_{1n}x_{n}=b_{1} \\
a_{21}x_{1}+ \cdots + a_{2n}x_{n} = b_{2} \\
\vdots \\
a_{m1}x_{1}+\cdots +a_{mn}x_{n}=b_{m}
\end{array}
\end{equation*}
Maka kita dapat menyatakan sistem ini dalam \textbf{bentuk matriks} \index{sistem persamaan!bentuk matriks} sebagai berikut.
\begin{equation*}
\leftB
\begin{array}{cccc}
a_{11} & a_{12} & \cdots & a_{1n} \\
a_{21} & a_{22} & \cdots & a_{2n} \\
\vdots & \vdots & \ddots & \vdots \\
a_{m1} & a_{m2} & \cdots & a_{mn}
\end{array}
\rightB
\leftB
\begin{array}{c}
x_{1} \\
x_{2} \\
\vdots \\
x_{n}
\end{array}
\rightB
=
\leftB
\begin{array}{c}
b_{1}\\
b_{2}\\
\vdots \\
b_{m}
\end{array}
\rightB
\end{equation*}
 
\end{definition}

Ekspresi $A\x=B$ juga dikenal sebagai \textbf{Bentuk Matriks} dari
sistem persamaan linear yang bersesuaian. \index{bentuk matriks $A\x=B$} Matriks
$A$ adalah matriks koefisien dari sistem, vektor
$\x$ adalah vektor kolom yang dibentuk dari variabel-variabel sistem,
dan akhirnya vektor $B$ adalah vektor kolom yang dibentuk dari konstanta-konstanta
sistem. Penting untuk diperhatikan bahwa setiap sistem persamaan linear
dapat ditulis dalam bentuk ini.

Perhatikan bahwa jika sistem persamaan homogen ditulis dalam bentuk matriks, bentuknya adalah
$A\x=0$, untuk vektor nol $0$.

Dari definisi ini dapat dilihat bahwa suatu vektor
\begin{equation*}
\x =
\leftB
\begin{array}{c}
x_{1} \\
x_{2} \\
\vdots \\
x_{n}
\end{array}
\rightB
\end{equation*}
akan memenuhi persamaan $A\x=B$
hanya jika entri-entri $x_{1}, x_{2}, \cdots, x_{n}$ dari vektor $\x$ merupakan solusi sistem semula.

Setelah menelaah cara mengalikan matriks dengan vektor, kini kita
ingin membahas kasus ketika kita mengalikan dua matriks dengan ukuran yang lebih
umum, meskipun ukurannya tetap harus sesuai sebagaimana akan kita
lihat. Sebagai contoh, dalam Contoh \ref{exa:vectormultbymatrix}, kita
mengalikan matriks $3 \times 4$ dengan vektor $4 \times 1$. Kita ingin
menyelidiki cara mengalikan matriks-matriks berukuran lain.

Kita belum memberikan syarat apa pun mengenai kapan perkalian matriks
dapat dilakukan! Untuk matriks $A$ dan $B$, agar hasil kali
$AB$ dapat dibentuk, banyaknya kolom $A$ harus sama dengan banyaknya baris
$B.$ Perhatikan hasil kali $AB$, dengan $A$ berukuran $m\times n$ dan $B$
berukuran $n \times p$. Maka hasil kali berdasarkan ukuran matriks
diberikan oleh
\begin{equation*}
(m\times\overset{\text{harus sama!}}{\widehat{n)\;(n}\times p})=m\times p
\end{equation*}

Perhatikan bahwa dua bilangan terluar memberikan ukuran hasil kali. Salah satu aturan terpenting tentang perkalian matriks adalah sebagai berikut.
Jika kedua bilangan di tengah tidak sama, matriks-matriks tersebut tidak dapat dikalikan!

Jika banyaknya kolom $A$ sama dengan banyaknya baris
$B$, kedua matriks disebut \textbf{konformabel}
\index{matriks!konformabel} dan hasil kali $AB$ diperoleh
\index{perkalian matriks} sebagai berikut.

\begin{definition}{Perkalian Dua Matriks}\label{def:multiplicationoftwomatrices}
 Misalkan $A$ adalah matriks $m\times n$
dan $B$ adalah matriks $n\times p$ berbentuk
\begin{equation*}
B=\leftB B_{1} \cdots  B_{p}\rightB
\end{equation*}
dengan $B_{1},...,B_{p}$ sebagai kolom-kolom $n\times 1$ dari $B$. Maka
matriks $m\times p$, yaitu $AB$, didefinisikan sebagai berikut:
\begin{equation*}
AB = A \leftB B_{1} \cdots  B_{p}\rightB =  \leftB (A B)_{1} \cdots  (AB)_{p}\rightB 
\end{equation*}
dengan $(AB)_{k}$ sebagai matriks $m\times 1$ atau vektor kolom yang
memberikan kolom $k$ dari $AB$.
\end{definition}

Perhatikan contoh berikut.

\begin{example}{Mengalikan Dua Matriks}\label{exa:multiplicationoftwomatrices}
Tentukan $AB$ jika memungkinkan.
\begin{equation*}
A = \leftB
\begin{array}{rrr}
1 & 2 & 1 \\
0 & 2 & 1
\end{array}
\rightB, B = \leftB
\begin{array}{rrr}
1 & 2 & 0 \\
0 & 3 & 1 \\
-2 & 1 & 1
\end{array}
\rightB
\end{equation*}
\end{example}

\begin{solution} Hal pertama yang perlu Anda periksa ketika menghitung suatu
hasil kali adalah apakah perkalian tersebut dapat dilakukan. Matriks pertama
berukuran $2\times 3$ dan matriks kedua berukuran $3\times 3$. Kedua
bilangan di dalam sama, sehingga $A$ dan $B$ merupakan matriks konformabel.
Berdasarkan pembahasan di atas, $AB$ akan berupa matriks $2\times 3$.
Definisi \ref{def:multiplicationoftwomatrices} memberi kita cara untuk
menghitung setiap kolom $AB$ sebagai berikut.
\begin{equation*}
\leftB \overset{
\text{Kolom pertama}}{\overbrace{\leftB
\begin{array}{rrr}
1 & 2 & 1 \\
0 & 2 & 1
\end{array}
\rightB \leftB
\begin{array}{r}
1 \\
0 \\
-2
\end{array}
\rightB }},\overset{\text{Kolom kedua}}{\overbrace{\leftB
\begin{array}{rrr}
1 & 2 & 1 \\
0 & 2 & 1
\end{array}
\rightB \leftB
\begin{array}{r}
2 \\
3 \\
1
\end{array}
\rightB }},\overset{\text{Kolom ketiga}}{\overbrace{\leftB
\begin{array}{rrr}
1 & 2 & 1 \\
0 & 2 & 1
\end{array}
\rightB \leftB
\begin{array}{r}
0 \\
1 \\
1
\end{array}
\rightB }}\rightB
\end{equation*}
Anda telah mengetahui cara mengalikan matriks dengan vektor, yaitu dengan menggunakan Definisi \ref{def:multiplicationvectormatrix} untuk
masing-masing dari ketiga kolom tersebut. Jadi,
\begin{equation*}
\leftB
\begin{array}{rrr}
1 & 2 & 1 \\
0 & 2 & 1
\end{array}
\rightB \leftB
\begin{array}{rrr}
1 & 2 & 0 \\
0 & 3 & 1 \\
-2 & 1 & 1
\end{array}
\rightB =\allowbreak \leftB
\begin{array}{rrr}
-1 & 9 & 3 \\
-2 & 7 & 3
\end{array}
\rightB 
\end{equation*}
\end{solution}

Karena vektor hanyalah matriks $ n \times 1$ atau $1 \times m$,
kita juga dapat mengalikan suatu vektor dengan vektor lain. 

\begin{example}{Perkalian Vektor dengan Vektor}\label{exa:vectormultiplication}
Kalikan, jika memungkinkan, $\leftB
\begin{array}{r}
1 \\
2 \\
1
\end{array}
\rightB \leftB
\begin{array}{rrrr}
1 & 2 & 1 & 0
\end{array}
\rightB .$ 
\end{example}

\begin{solution} Dalam hal ini, kita mengalikan matriks berukuran $3 \times 1$ dengan matriks berukuran $1 \times
4.$ Kedua bilangan di dalam sama, sehingga hasil kali tersebut terdefinisi. Perhatikan bahwa hasil kalinya berupa matriks berukuran $3 \times 4$.
Dengan menggunakan Definisi \ref{def:multiplicationoftwomatrices}, kita dapat menghitung hasil kali ini sebagai berikut $\:$
\begin{equation*}
\leftB
\begin{array}{r}
1 \\
2 \\
1
\end{array}
\rightB \leftB
\begin{array}{rrrr}
1 & 2 & 1 & 0
\end{array}
\rightB = 
\leftB \overset{
\text{Kolom pertama}}{\overbrace{\leftB
\begin{array}{r}
1 \\
2 \\
1
\end{array}
\rightB \leftB
\begin{array}{r}
1
\end{array}
\rightB }},\overset{\text{Kolom kedua}}{\overbrace{\leftB
\begin{array}{r}
1 \\
2\\
1
\end{array}
\rightB \leftB
\begin{array}{r}
2 
\end{array}
\rightB }},\overset{\text{Kolom ketiga}}{\overbrace{\leftB
\begin{array}{r}
1 \\
2 \\
1
\end{array}
\rightB \leftB
\begin{array}{r}
1
\end{array}
\rightB }}, \overset {\text{Kolom keempat}}{\overbrace{\leftB
\begin{array}{r}
1\\
2\\
1
\end{array}
\rightB \leftB
\begin{array}{r}
0
\end{array}
\rightB}}
\rightB
\end{equation*}

Anda dapat menggunakan Definisi \ref{def:multiplicationvectormatrix} untuk memeriksa bahwa hasil kali ini adalah
\begin{equation*}
\leftB
\begin{array}{cccc}
1 & 2 & 1 & 0 \\
2 & 4 & 2 & 0 \\
1 & 2 & 1 & 0
\end{array}
\rightB
\end{equation*}
\end{solution}

\begin{example}{Perkalian yang Tidak Terdefinisi}\label{exa:undefinedmatrixmultiplication}
Tentukan $BA$ jika memungkinkan.
\begin{equation*}
B = \leftB
\begin{array}{ccc}
1 & 2 & 0 \\
0 & 3 & 1 \\
-2 & 1 & 1
\end{array}
\rightB,  A = \leftB
\begin{array}{ccc}
1 & 2 & 1 \\
0 & 2 & 1
\end{array}
\rightB
\end{equation*}
\end{example}

\begin{solution} Pertama, periksa apakah perkalian dapat dilakukan. Hasil kali ini berbentuk $\left( 3\times 3\right)
\left( 2\times 3\right) .$ Kedua bilangan di dalam tidak sama sehingga perkalian ini tidak dapat
dilakukan. 
\end{solution}

Dalam hal ini, kita mengatakan bahwa perkalian tersebut tidak terdefinisi.
Perhatikan bahwa matriks-matriks ini sama dengan yang kita gunakan dalam Contoh \ref{exa:multiplicationoftwomatrices}.
Dalam contoh ini, kita mencoba menghitung $BA$, bukan $AB$. Hal ini menunjukkan sifat lain
dari perkalian matriks. Meskipun hasil kali $AB$ mungkin terdefinisi, kita tidak dapat mengasumsikan bahwa
hasil kali $BA$ dapat dilakukan. Oleh karena itu, penting untuk selalu memeriksa bahwa hasil kali terdefinisi
sebelum melakukan perhitungan apa pun.

Sebelumnya, kita mendefinisikan matriks nol $0$ sebagai matriks (dengan
ukuran yang sesuai) yang seluruh entrinya nol. Perhatikan
contoh perkalian dengan matriks nol berikut.

\begin{example}{Perkalian dengan Matriks Nol}\label{exa:multbyzeromatrix}
Hitung hasil kali $A0$ untuk matriks
\begin{equation*}
A=
\leftB
\begin{array}{rr}
1 & 2 \\
3 & 4
\end{array}
\rightB
\end{equation*}
dan matriks nol $2 \times 2$ yang diberikan oleh
\begin{equation*}
0=
\leftB
\begin{array}{rr}
0 & 0 \\
0 & 0
\end{array}
\rightB
\end{equation*}
\end{example}

\begin{solution} 
Dalam hasil kali ini, kita menghitung
\begin{equation*}
\leftB
\begin{array}{rr}
1 & 2 \\
3 & 4
\end{array}
\rightB
\leftB
\begin{array}{rr}
0 & 0 \\
0 & 0
\end{array}
\rightB
=
\leftB
\begin{array}{rr}
0 & 0 \\
0 & 0
\end{array}
\rightB
\end{equation*}

Jadi, $A0=0$. 
\end{solution}

Perhatikan bahwa kita juga dapat mengalikan $A$ dengan vektor nol $2 \times 1 $ yang diberikan oleh $\leftB
\begin{array}{r}
0 \\
0
\end{array}
\rightB$.
Hasilnya adalah vektor nol $2 \times 1$.
Oleh karena itu, selalu berlaku $A0=0$ untuk matriks atau vektor nol dengan ukuran yang sesuai.
